%% kapa: exam=1 round=31 year=2020 period=1 form=A subject=civil_law subject_ko=민법 seq=23 number=23
%% source: Q-Net 감정평가사 기출문제 (cst00309) artlSeq=5207219 file=extracted/1008/5207219/A형/A형/2020년도 제31회 감정평가사 1차 시험_1교시_A형.pdf
%% conversion: kordoc 4.24.51 -> tools/gen1.py (auto); review=proofread-sonnet; flags=none
\begin{question}{23}
물권적 청구권에 관한 설명으로 옳지 않은 것은? (다툼이 있으면 판례에 따름)
\begin{choices}
  \item 물권적 청구권은 물권과 분리하여 양도하지 못한다.
  \item 물권적 청구권을 보전하기 위하여 가등기를 할 수 있다.
  \item 미등기건물을 매수한 사람은 소유권이전등기를 갖출 때까지 그 건물의 불법점유자에게 직접 자신의 소유권에 기하여 인도를 청구하지 못한다.
  \item 토지소유자는 권원 없이 그의 토지에 건물을 신축ㆍ소유한 사람으로부터 건물을 매수하여 그 권리의 범위에서 점유하는 사람에게 건물의 철거를 청구할 수 있다.
  \item 소유권에 기한 말소등기청구권은 소멸시효의 적용을 받지 않는다.
\end{choices}
\end{question}
\answer{2}
